\documentclass{article}

\usepackage[T1]{fontenc}

\usepackage{booktabs,graphicx,siunitx}

\begin{document}

\section{Fisher's exact}

\begin{tabular}{llll}
\toprule
Row \textbackslash{} Column & No & Yes & Total \\
\midrule
Female & 44 & 76 & 120 \\
Male & 43 & 77 & 120 \\
Total & 87 & 153 & 240 \\
\bottomrule
\end{tabular}

\begin{tabular}{llll}
\toprule
p (exact) & Odds ratio (cond. MLE) & 95\% CI & Cram\'{e}r's V [95\% CI] \\
\midrule
1.000 & 1.04 & [0.59, 1.82] & --- \\
\bottomrule
\end{tabular}

\includegraphics{figures/01_fishers-exact/figure_bar.png}

An exact alternative to chi-square for small samples. A p below alpha means the two categories are associated; for 2$\times$2 tables the odds ratio quantifies the association --- this is fisher.test()'s conditional MLE (the non-central hypergeometric estimate), not the sample cross-product. OR$>$1 means the outcome is more likely in the first row/group, OR$<$1 less likely, OR=1 no association; check the row/column order before reading direction, since the OR reflects the table's orientation. For tables larger than 2$\times$2 the odds ratio is undefined, so Cram\'{e}r's V (0--1) reports the strength of association instead. Your significance threshold ($\alpha$) is 0.05.

A Fisher's exact test of gender by passed\_course gave p = 1.000 (OR=1.04, 95\% CI [0.59, 1.82]).

Statistical basis: Fisher's exact test (conditional MLE odds ratio for 2x2 tables) (Fisher, R. A. (1922). "On the interpretation of $\chi$\textsuperscript{2} from contingency tables, and the calculation of P." Journal of the Royal Statistical Society, 85(1), 87-94.); Cram\'{e}r's V effect size (larger-than-2$\times$2 tables) (Ben-Shachar MS, Lüdecke D, Makowski D (2020). "effectsize: Estimation of Effect Size Indices and Standardized Parameters." Journal of Open Source Software, 5(56), 2815.).

\end{document}